\subsection{Witt designs and their actual
groups}\label{witt-designs-and-their-actual-groups}

The program \texttt{phase186\_witt\_derived\_check.g} constructs the five native
Mathieu designs from group orbits. For M₁₁ and M₁₂ the orbits on five-
and six-subsets have sizes 66,396 and132,792, respectively. Thus exactly
one orbit can be the block set of the required Steiner system. For each
large Mathieu group, the pointwise stabilizer of a fixed ordered t-tuple
has two remaining point orbits, of sizes 3 and16. The unique block
through the tuple must consist of the tuple and the size-three orbit.
This proves uniqueness of the native invariant block system before the
actual Steiner condition is checked.

The program then checks that each t-subset occurs in exactly one block.
The full design automorphism group is computed afresh from the incidence
structure, its generators are checked to preserve the complete block
set, and its relation to the original Mathieu subgroup is tested as an
actual permutation-group inclusion or equality.

{\def\LTcaptype{none} % do not increment counter
\begin{longtable}[]{@{}
  >{\raggedright\arraybackslash}p{(\linewidth - 8\tabcolsep) * \real{0.1579}}
  >{\raggedleft\arraybackslash}p{(\linewidth - 8\tabcolsep) * \real{0.2105}}
  >{\raggedleft\arraybackslash}p{(\linewidth - 8\tabcolsep) * \real{0.2105}}
  >{\raggedleft\arraybackslash}p{(\linewidth - 8\tabcolsep) * \real{0.2105}}
  >{\raggedleft\arraybackslash}p{(\linewidth - 8\tabcolsep) * \real{0.2105}}@{}}
\toprule\noalign{}
\begin{minipage}[b]{\linewidth}\raggedright
Native model
\end{minipage} & \begin{minipage}[b]{\linewidth}\raggedleft
Blocks
\end{minipage} & \begin{minipage}[b]{\linewidth}\raggedleft
t-subsets checked
\end{minipage} & \begin{minipage}[b]{\linewidth}\raggedleft
Native group order
\end{minipage} & \begin{minipage}[b]{\linewidth}\raggedleft
Full design group order
\end{minipage} \\
\midrule\noalign{}
\endhead
\bottomrule\noalign{}
\endlastfoot
S(4,5,11) & 66 & 330 & 7,920 & 7,920 \\
S(5,6,12) & 132 & 792 & 95,040 & 95,040 \\
S(3,6,22) & 77 & 1,540 & 443,520 & 887,040 \\
S(4,7,23) & 253 & 8,855 & 10,200,960 & 10,200,960 \\
S(5,8,24) & 759 & 42,504 & 244,823,040 & 244,823,040 \\
\end{longtable}
}

For degrees11,12,23 and24 the full group equals the native group, not
merely a group of the same order. At degree22 the native M₂₂ is a
subgroup of index two. This supplies the full-group bounds in the
small-Witt exclusions: the chamber counts23,760 and285,120 exceed the
respective full group orders by a factor of three.

\subsubsection{Explicit maps to derived projective
planes}\label{explicit-maps-to-derived-projective-planes}

The target plane is constructed directly from GF(4)³. Points are nonzero
vectors normalized so that their first nonzero coordinate is1. Lines are
the kernels of the21 similarly normalized nonzero linear forms. Hence
the model is PG(2,4) by its coordinate definition, without invoking
uniqueness of projective planes of order 4.

For each large Witt design, remove the fixed (t−2)-set
R=\{1,\ldots,t−2\} from every block containing R. The program finds an
isomorphism of the resulting two-colored incidence graph with the
coordinate plane. It then discards the graph-search conclusion as a
black box and verifies the explicit point permutation: the mapped block
set equals the coordinate line set, and all441 point/line incidences and
nonincidences agree. Native t-transitivity transports this conclusion to
every (t−2)-fold derivative.

The saved JSON witness records all original blocks, all derived blocks,
the21 coordinate vectors, the coordinate lines, and each explicit point
map. The independent \texttt{phase186\_witt\_validate.py} recomputes GF(4)
arithmetic as coefficient-bit arithmetic modulo z²+z+1, reconstructs all
coordinate lines, reconstructs each derivative from its original blocks,
and checks each map. It also checks every Steiner condition and the
design-preservation action of every saved group generator. The validator
does not claim an independent proof that the saved automorphism
generators exhaust the full group.

The file formal/WittDerivedPlanes.lean verifies explicit point and line
bijections for all three displayed derivatives and proves that all
incidences agree with the literal GF(4) dot-product equation. It reuses
the frozen F4 arithmetic definitions and verified field-table laws in
DeletedLineFourResidue. All nine retained axiom reports use only propext,
Classical.choice and Quot.sound. There is no assumed
incidence-isomorphism axiom, sorry, admit or native\_decide. Two
deliberately altered coordinate-incidence maps are rejected by Lean in
the fresh replay. The generator binds the sorted GAP line-support
convention to the separately ordered coordinate linear forms; the
resulting source is checked byte-for-byte against the reviewed Lean
file.

This formal certificate starts from the displayed finite derivative
tables. Its identification of those tables as derivatives of the
original Witt blocks remains outside Lean, checked independently by GAP
and Python. The all-rank residue recursion, full design automorphism
groups, and stabilizer structures also remain outside this Lean
certificate.

\subsubsection{The displayed sizes and local orders}
Theorem S.3 gives the following object sizes, in the order ordinary
subset types, marked objects, blocks, copied points. The panel sizes
are listed in the same order; the corresponding local orders are one
less. This states the numerical Figure 3 data independently of the
group identifications.
\begin{center}\small
\begin{tabular}{lll}
\toprule
Design & Object sizes & Panel sizes\\
\midrule
$S(2,5,21)$ & $105,21,21$ & $5,4,4$\\
$S(3,6,22)$ & $22,1155,77,22$ & $2,5,4,4$\\
$S(4,7,23)$ & $23,253,8855,253,23$ & $2,2,5,4,4$\\
$S(5,8,24)$ & $24,276,2024,53130,759,24$ & $2,2,2,5,4,4$\\
\bottomrule
\end{tabular}
\end{center}
Indeed the block and marked counts are respectively
$\binom vt/\binom kt$ and that number times $\binom k{t-1}$;
the ordinary counts are $\binom vj$. All four rows have $r=5$.
The coordinate-plane row is the $q=4$ case of B.ProjectivePlane;
its rank-two residues are affine planes and a generalized digon.
For the three Witt rows the common residue decomposition S.4 and the
explicit derived-plane maps give the same terminal residue data,
with the Boolean intervals supplying the ordinary chains. Thus the
geometric meaning of the diagram is not inferred from the numerical
counts alone.

\subsubsection{Block and marked-object
stabilizers}\label{block-and-marked-object-stabilizers}

Let G be the native group, B a block of size k, and H=G\_B. For k=6,7,8,
restriction to B gives an exact sequence \[
1\longrightarrow N\longrightarrow H\longrightarrow A_k\longrightarrow1.
\] The computation proves that N is elementary abelian of order 16, the
image is the actual alternating permutation group on B, and a complement
C exists with N∩C=1 and NC=H. Thus H=2⁴:A\_k, with the action supplied
by the saved permutation realization. An order factorization alone is
not used to infer splitting.

A marked object is (T,B∖T), where \textbar T\textbar=t−1. Its stabilizer
M is exactly H\_T. Since N fixes every point of B, it fixes T, and \[
M=N\rtimes C_T,\qquad
C_T\cong (S_{t-1}\times S_{k-t+1})\cap A_k.
\tag{W1}
\] The program verifies the actual image equality, the complement
intersection and the product order in (W1). These give all block and
marked labels in the three Mathieu rows of Figure3:

{\def\LTcaptype{none} % do not increment counter
\begin{longtable}[]{@{}llll@{}}
\toprule\noalign{}
Group & Block stabilizer & Marked stabilizer & Orders \\
\midrule\noalign{}
\endhead
\bottomrule\noalign{}
\endlastfoot
M₂₂ & 2⁴:A₆ & 2⁴:S₄ & 5,760;384 \\
M₂₃ & 2⁴:A₇ & 2⁴:((C₃×A₄):C₂) & 40,320;1,152 \\
M₂₄ & 2⁴:A₈ & 2⁴:((A₄×A₄):C₂) & 322,560;4,608 \\
\end{longtable}
}

For the latter two descriptions, take the even permutations in the two
disjoint factors. Their product has index two in the parity-matched
subgroup; a pair of transpositions, one in each factor, gives a
complementary involution. This proves the displayed semidirect
descriptions directly.

The source writes the M₂₂ marked group as 2²·2⁴:S₃. To bind that
notation as well, the program computes Q=O₂(M), proves
\textbar Q\textbar=64, Z(Q)≅2² and Q/Z(Q)≅2⁴, proves M/Q≅S₃, and finds a
complement to Q. Thus both descriptions refer to the same explicitly
constructed group. The dot in the source's middle extension is retained;
no unwarranted splitting of Q over its center is asserted.

For the first row, G=PSL(3,4) acts on the coordinate plane. Applying the
identical block calculation to a line of size 5 gives the line
stabilizer2⁴:A₅, and fixing a point on that line gives the marked
stabilizer2⁴:A₄. Projective duality, which normalizes PSL(3,4), gives
the same2⁴:A₅ point stabilizer. Their orders are 960 and192.

\subsubsection{Ordinary-subset
stabilizers}\label{ordinary-subset-stabilizers}

For an ordinary j-subset R, the set stabilizer G\_R acts on the j-fold
derived design. The computation verifies faithfulness of this action,
finds an explicit incidence map to the appropriate smaller native Witt
design or coordinate projective plane, and conjugates the actual
residual permutation group by that map.

The pointwise stabilizer of R becomes precisely the smaller native
Mathieu group, or precisely the group generated by the projective action
of SL(3,4) when j=t−2. This is a group equality in the target
permutation action. The set stabilizer gives:

{\def\LTcaptype{none} % do not increment counter
\begin{longtable}[]{@{}lrl@{}}
\toprule\noalign{}
Native group & Ordinary subset size & Actual identified stabilizer \\
\midrule\noalign{}
\endhead
\bottomrule\noalign{}
\endlastfoot
M₂₂ & 1 & PSL(3,4), order 20,160 \\
M₂₃ & 1 & M₂₂, order 443,520 \\
M₂₃ & 2 & PSL(3,4):C₂, order 40,320 \\
M₂₄ & 1 & M₂₃, order 10,200,960 \\
M₂₄ & 2 & full Aut(S(3,6,22)), order 887,040 \\
M₂₄ & 3 & PΓL(3,4)=PSL(3,4):S₃, order 120,960 \\
\end{longtable}
}

The source labels the penultimate group M₂₂:2. The current full-group
computation gives the actual degree22 overgroup; identifying the colon
as a split extension additionally requires a complementary involution,
which is checked in the finite runner. The projective index-two case is
also checked to split. For the final case the program builds PΓL(3,4)
explicitly from GL(3,4) acting on normalized vectors together with the
coordinate Frobenius map, and verifies equality with the transported
stabilizer. Its quotient over PSL(3,4) is C₃:C₂: determinant modulo
cubes gives C₃, and Frobenius inverts it. A diagonal matrix of order 3
and the Frobenius map supply a complement, proving the displayed S₃
extension.

Copied-point stabilizers are ordinary-point stabilizers. This proves the
Figure 3 isomorphism labels, with explicit actions rather than
recognition from orders.
